% Pursuing Stacks (cotes 134-2 to 134-8): archive cote 134-6, site page 36 (PDF page 37), typescript page 354.
% typescript. Draft compared: ps5.tex:3102-3166.
% AI-transcribed by Claude Opus 5.5 (claude-opus-5-5, adaptive thinking, effort high), Message Batches API, batch ps3 (msgbatch_01XdFmuVVBAcgc7EFxdtCqsj), 2026-11-12; not yet checked by eye.
\scanpage{134-6/36}
\mspage{354}

We were faced yesterday with three different homotopy structures
$h,h',h''$ on a category $A\hat{}$, for a given small category $A$, which make
a sense for any $A$, and which in case $A$ is a contractor all coincide.
The exact relationship between these structures in more general cases
has remained somewhat confused, and in order to dispell the \unsure{fesulting|resulting}
feeling of uneasiness, I took finally the trouble \inserted{today} to write it out with
some care. One of these structures, $h'$, depend\hand{s} on the choice of a
basic localizer $\underline{W}$ , whereas the two others do'nt... \xout{\illegible{1}}

12.7.

97. I was interrupted in my notes by visiting friends arriving in
close succession - then since yesterday I have been busy mainly with
letter writing. Now I am ready to take up the thread where I left it
- namely some afterthoughts to the reflections of section 95 on con-%
tractors.

First an afterthought to the afterthoughts ! I had introduced
the name ``\uline{vertical} topos'' for a topos admitting only the two trivial
open subtopoi (page 352), whereas the stronger property that every
``non-empty'' sheaf has a section remained unnamed (which is no real
drawback as long as we are restricting to topoi\hand{/}of the type $A\hat{}$ , where
\inserted{indeed} the two notions coïncide). Now the latter property can be viewed as
the combination of two others, one being verticality, and the other
the property that every sheaf $F$ \uline{such that} $F \to e$ \uline{be epimorphic},
should admit a section. It is this la\over{tter}{st} property which does merit
to `` be viewed as the strongest conveivable asphericity \inserted{property} for a topos '',
as I commented on it last friday (p.352). After I had written this
down as a kind of selfevidence, \xout{I had} a doubt turned up though and I
qualified the comment by adding ``as far as $H^1$ goes at any rate, as
the $H^1$ with \xout{values} \inserted[typed]{coefficients} in \uline{any} group object will be zero'' . I did'nt pau\xout{s}
se then to see if the doubt was founded - quite evidently it is'nt\hand{,}
\inserted{except for $H^0$,} as it is clear by the usual shift argument, using embedding of \xout{abelia}
an abelian sheaf into an injective one, that if \xout{the \unsure{$H^1$}} for given
$k$ (here $k=1$) $H^k(X,F) = 0$ for any abelian sheaf $F$, then the same
holds for $H^n$ with any $n \hand{\geqslant} k$ - i.e. the \inserted{global} \unsure{conomological|cohomological} dimension of $X$
is $\hand{<} k$ . This implies that \xout{if} any small vertical category (a forti-%
ori any \xout{contrx} \uline{precontractor}) is aspheric\inserted{, provided it is \unsure{connected}}, indeed its cohomology in-%
\margin{0-connected}
variants with values in \uline{any} sheaf of coefficients (not necessarily
commutative as far as $H^1$ goes) are trivial.

\margin{36}
